\subsection{Relative locality of the charged and neutral fields}

\begin{theorem}[Relative locality survives the graph limit]
\label{thm:relativeLocality}
The limiting edge-projected charged field is relatively local with
respect to the neutral $D_5\oplus A_3$ current net.
\end{theorem}

\begin{proof}
For every neutral root $\alpha$, the frozen generator satisfies
\[
 \ip{\lambda}{\alpha}\in\Z.
\tag{5.26}\label{eq:integralpairing}
\]
On the displayed simple-root basis the eight pairings are
$[0,0,0,0,1,1,0,0]$; integrality for the full root lattice follows by
linearity.  Therefore the finite vertex exchange phase is
$e^{2\pi i\ip{\lambda}{\alpha}}=1$.  If $f$ and $g$ have disjoint
supports, the finite collar consequently satisfies
\[
 [\Psi_{\lambda,N}(f),W_{\alpha,N}(g)]\xi=0
\]
on the finite-energy core, where $W_{\alpha,N}(g)$ is a bounded neutral
Weyl generator.

Transport to the common Hilbert space with \eqref{eq:Fockisometry}.
Theorem~\ref{thm:TELa} gives graph convergence of the charged field and
its adjoint; Osborne--Stottmeister gives strong convergence of the
bounded neutral Weyl operators.  For a common-core vector $\xi$,
\[
\begin{aligned}
 W_\alpha(g)\Psi_\lambda(f)\xi
 &=\lim_NW_{\alpha,N}(g)\Psi_{\lambda,N}(f)\xi\\
 &=\lim_N\Psi_{\lambda,N}(f)W_{\alpha,N}(g)\xi
 =\Psi_\lambda(f)W_\alpha(g)\xi .
\end{aligned}
\]
The uniform graph bound justifies both products.  Closability extends
the equality from the core, and the Weyl generators strongly generate
the neutral local algebra.
\end{proof}

\subsection{Outerness of the limiting four-grade action}

\begin{theorem}[Outerness from the connected simple-current Q-system]
\label{thm:outerness}
On every limiting local factor generated by the neutral algebra and the
four charged grades, every nonidentity element of the $\Z_4$ gauge
action is outer.
\end{theorem}

\begin{proof}
Let $\iota:\A(I)\hookrightarrow\B^{\rm gen}(I)$ be the local inclusion.
Its canonical endomorphism has the multiplicity-free decomposition
\[
 \overline\iota\iota
 \simeq\id\oplus\rho_\lambda\oplus\rho_{2\lambda}
              \oplus\rho_{3\lambda}.
\tag{5.27}\label{eq:canonicalendo}
\]
The four sectors are inequivalent: their exact grades are distinct and
$[\lambda]$ has order four.  Only the first summand is the vacuum, so
\[
 \A(I)'\cap\B^{\rm gen}(I)
 \cong\operatorname{Hom}(\id,\overline\iota\iota)=\C.
\tag{5.28}\label{eq:irreducibleinclusion}
\]
Suppose $\gamma_g=\operatorname{Ad}u$ for a nontrivial $g\in\Z_4$ and
$u\in\B^{\rm gen}(I)$.  Since $\gamma_g$ fixes $\A(I)$ pointwise,
$u\in\A(I)'\cap\B^{\rm gen}(I)=\C$, so $\gamma_g=\id$.  This contradicts
faithfulness: on the grade-one field it acts by the nontrivial character
$\gamma_g(\Psi_\lambda)=\chi_1(g)\Psi_\lambda$.  Hence $\gamma_g$ is
outer.

There is no finite-dimensional outerness witness: every automorphism of
$M_{256}$ is inner, and
$U=\operatorname{diag}(1_{64},i1_{64},-1_{64},-i1_{64})$ is a
charge-zero implementing unitary.  Outerness is a local-factor limit
statement derived from \eqref{eq:irreducibleinclusion}, not inherited
from finite matrix algebras.
\end{proof}

\subsection{From fields to local algebras}

For an interval $I$, define the field-generated finite algebra
\[
 \B_N^{\rm gen}(I)
 =W^*\!\left(
 \A_N(I),\
 \overline{\Psi_{k\lambda,N}(f)}:
 k\in\Z_4,\ \supp f\subset I
 \right).
\tag{5.29}\label{eq:generatedfinite}
\]
Define $\B^{\rm gen}(I)$ similarly from the limiting fields.

\begin{lemma}[Graph convergence of a complete generating family]
\label{lem:algconv}
Assume L1, the graph convergence \eqref{eq:graphconv} for a
support-preserving countable dense family of test functions, strong
locality of the limiting fields, and the exhaustion condition
\textup{(ALG-EXH)} from the final box.  Then
\[
 \B_N^{\rm gen}(I)\longrightarrow\B^{\rm gen}(I)
\tag{5.30}\label{eq:algconv}
\]
in the standard scaling-algebra sense: bounded functions of the
generating fields converge strongly, the strong lower and upper limits
coincide, and no additional local scaling operator remains.
\end{lemma}

\begin{proof}
Graph convergence of a closed operator and its adjoint implies strong
resolvent convergence for its hermitian real and imaginary parts on the
common invariant core.  Bounded continuous functions of those parts
therefore converge strongly.  Their spectral projections generate the
same von Neumann algebra as the affiliated smeared fields.  Thus every
element of the limiting generated algebra belongs to the strong lower
limit.  Conversely, \textup{(ALG-EXH)} says exactly that every bounded local
scaling sequence is approximable by bounded words in this generating
family; its weak cluster points lie in $\B^{\rm gen}(I)$.  Hence the
upper limit is contained in the lower limit.
\end{proof}

\begin{proposition}[Sharp Buchholz--Verch exhaustion reduction]
\label{prop:BVexhaustion}
Let $\underline{\B}^{\rm gen}(I)$ denote the scaling subalgebra generated
by the neutral Weyl operators and the four graph-convergent charged
grades.  Equality with the full microscopic Buchholz--Verch scaling
algebra follows if the following two estimates hold:
\[
\begin{split}
 &\lim_{E\to\infty}\ \sup_{N\ge16}\
 \sup_{\substack{x\in\B_N(I)\\\norm{x}\le1}}
 \norm{(1-P^{\rm gen}_{E,N})e^{-\beta L_{0,N}}x\Omega_N}=0
 \quad(\beta>0),\\
 &\lim_{N\to\infty}\ 
 \inf_{y_N\in\B_N^{\rm gen}(I)}
 \left(
 \norm{(x_N-y_N)P_{E,N}}
 +\norm{(x_N^*-y_N^*)P_{E,N}}
 \right)=0
\end{split}
\tag{5.31}\label{eq:EXHsharp}
\]
for every finite $E$ and every norm-bounded, translation-continuous
local scaling sequence $(x_N)$.
\end{proposition}

\begin{proof}
The first line is uniform phase-space tightness: after Euclidean energy
damping, vectors created by the microscopic unit ball are uniformly
approximated by the finite-energy neutral-plus-four-grade subspace.  The
second line says that no bounded operator sequence has additional
matrix elements on any finite-energy sector, for the operator and its
adjoint.  It gives equality of the strong lower limits on the common
finite-energy core.  Phase-space tightness then excludes a weak cluster
point supported only at energies escaping to infinity, giving equality
of the upper limit as well.  This is the standard scaling-algebra
compactness mechanism of Buchholz--Verch \cite{BuchholzVerch}, written
here as the two concrete estimates the collar must satisfy.
\end{proof}

\begin{remark}[Verified and missing Buchholz--Verch hypotheses]
The frozen lattice system has isotony, exact finite-range locality,
uniform norm bounds for scaling families, discrete translation
covariance, and a translation-covariant neutral scaling limit.
Theorems~\ref{thm:UGF}, \ref{thm:TELa} and
\ref{thm:relativeLocality} add uniform charged energy control, graph
convergence and the lower inclusion
$\underline{\B}^{\rm gen}(I)\subset\underline{\B}(I)$.  What is not
proved is the uniform microscopic phase-space estimate and the second,
upper-inclusion line of \eqref{eq:EXHsharp}.  These are precisely the
``no extra ultraviolet operators'' content; no arithmetic or finite
sector census can establish them.
\end{remark}

\begin{remark}[Binding ALG2 correction, 2026-09-05]
The former claim that Fourier/Vandermonde interpolation automatically proves
the second line of \eqref{eq:EXHsharp} for
$\B_N^{\rm gen}(I)$ is withdrawn.  The exact one-mode countermodel in
\texttt{v1022\_telb\_round3\_certificates.py} has full CAR algebra $M_2$,
whereas its even/code-generated subalgebra is diagonal; the annihilator is at
distance at least $1$ from that subalgebra.  Vandermonde inversion reaches a
full-CAR compression, but does not identify the microscopic generated algebra.
Consequently the upper inclusion requires an independent FE-GEN
compression-surjectivity theorem, or a precise gauge/code-projected
microscopic-algebra argument.  The positive phase-space/ALG1 reading is also
conditional on the generated projection containing the microscopic
low-energy sectors.
\end{remark}

\begin{theorem}[Exact reduction of L2]\label{thm:L2reduction}
Assume the exhaustion estimates \eqref{eq:EXHsharp}.  Then L2 of
Theorem~\ref{thm:master} holds.
\end{theorem}

\begin{proof}
Lemma~\ref{lem:telescoping} gives the nonzero charged field after fixing
one coherent cocycle/phase normalization.
Theorem~\ref{thm:relativeLocality} and the integer spin $h_\lambda=1$
make the order-four Q-system commutative.
The neutral stress tensor converges by Theorem~\ref{thm:OS}; covariance
of the charged primary fields then supplies conformal covariance of the
generated extension.  Lemma~\ref{lem:algconv} gives local-algebra
convergence, using Proposition~\ref{prop:BVexhaustion}.
Theorem~\ref{thm:outerness} supplies the outer local action.
Nonzeroness follows from the normalized vacuum-to-charge
matrix element fixed to one at every scale and preserved by the graph
limit.
\end{proof}

\section{\texorpdfstring{L3: finite-group $C^*$ crossed products commute with the limit}{L3: finite-group crossed products commute with the limit}}
\label{sec:L3}

\subsection{\texorpdfstring{The $C^*$-inductive theorem}{The C-star inductive theorem}}

\begin{theorem}[Finite-group crossed-product/limit exchange]
\label{thm:crossedlimit}
Let $G$ be finite and let
\[
 A_1\xrightarrow{\iota_1}A_2\xrightarrow{\iota_2}\cdots
\]
be an inductive system of unital $C^*$-algebras with injective maps.
Suppose $\alpha^{(N)}:G\curvearrowright A_N$ are actions and
\[
 \iota_N\alpha_g^{(N)}
 =\alpha_g^{(N+1)}\iota_N .
\tag{6.1}\label{eq:equivariance}
\]
Let $A=\varinjlim A_N$ and $\alpha=\varinjlim\alpha^{(N)}$.  Then the
canonical map is a $*$-isomorphism
\[
 \varinjlim(A_N\rtimes_{\alpha^{(N)}}G)
 \cong A\rtimes_\alpha G.
\tag{6.2}\label{eq:crossediso}
\]
For the coefficient expectations
\[
 F_N\!\left(\sum_{g\in G}a_gu_g\right)=a_e
\tag{6.3}\label{eq:coeffE}
\]
and fixed-point expectations
\[
 E_N(a)=\frac1{|G|}\sum_{g\in G}\alpha_g^{(N)}(a),
\tag{6.4}\label{eq:fixedE}
\]
the diagrams commute exactly, and their limits are the canonical
expectations $F$ and $E$.  They are faithful.  Moreover,
\[
 E_N(x)\ge\frac1{|G|}x\quad(x\ge0),
\qquad
 E(x)\ge\frac1{|G|}x\quad(x\ge0).
\tag{6.5}\label{eq:PP}
\]
Thus the uniform Pimsner--Popa index bound survives the limit.  If each
action is saturated/outer on factors, the index is exactly $|G|$.
\end{theorem}

\begin{proof}
Equivariance gives a homomorphism
\[
 j_N:A_N\rtimes G\longrightarrow A_{N+1}\rtimes G,\qquad
 j_N(a)=\iota_N(a),\quad j_N(u_g)=u_g.
\tag{6.6}\label{eq:crossmap}
\]
It is injective.  To see this without appealing to exactness, choose a
faithful representation $\pi_{N+1}$ of $A_{N+1}$ and use the regular
covariant representation on $\ell^2(G)\otimes\Hh$.  Its restriction to
$A_N\rtimes G$ is faithful because finite groups are amenable, so full
and reduced crossed products coincide.  Hence the inductive limit on the
left of \eqref{eq:crossediso} is defined by injective maps.

Let $A_{\rm alg}=\bigcup_N\iota_{N,\infty}(A_N)$, dense in $A$.
Every algebraic crossed-product element
$x=\sum_{g\in G}a_gu_g$ with $a_g\in A_{\rm alg}$ lies at one common
finite stage because $G$ is finite.  Therefore the union of the images
of $A_N\rtimes_{\alpha^{(N)}}G$ is the algebraic crossed product
$A_{\rm alg}\rtimes G$, which is dense in $A\rtimes_\alpha G$.
The regular representations used above show that the norm of such an
$x$ is the same whether computed at its finite stage, in the inductive
limit, or in $A\rtimes_\alpha G$.  The dense isometric homomorphism
therefore extends to \eqref{eq:crossediso}.

On algebraic elements,
\[
 F_{N+1}j_N\!\left(\sum_ga_gu_g\right)
 =\iota_N(a_e)=\iota_NF_N\!\left(\sum_ga_gu_g\right),
\]
and \eqref{eq:equivariance} gives
$E_{N+1}\iota_N=\iota_NE_N$.  Hence both expectations descend to the
limit and agree there with \eqref{eq:coeffE} and
\eqref{eq:fixedE}.  Faithfulness of $F$ follows from the faithful regular
representation.  Faithfulness of $E$ follows because a sum of positive
operators can vanish only when every summand vanishes.

Finally, for $x\ge0$,
\[
 E_N(x)=\frac1{|G|}
 \left(x+\sum_{g\ne e}\alpha_g^{(N)}(x)\right)
 \ge\frac1{|G|}x.
\]
The same direct formula proves the limit inequality.  For outer finite
group actions on factors, the fixed-point index theorem gives equality
of the index with $|G|$.
\end{proof}

\subsection{State preservation and the charged grades}

\begin{corollary}[State-preserving expectations commute with the limit]
\label{cor:stateE}
Let $\varphi_N$ be compatible $\alpha^{(N)}$-invariant states.  Then
$\varphi_N\circ E_N=\varphi_N$, the limit state $\varphi$ satisfies
$\varphi\circ E=\varphi$, and the GNS extensions of $E_N$ converge on
the dense inductive union to the GNS extension of $E$.
\end{corollary}

\begin{proof}
Invariance gives
\[
 \varphi_N(E_N(a))
 =|G|^{-1}\sum_g\varphi_N(\alpha_g^{(N)}(a))
 =\varphi_N(a).
\]
Compatibility permits passage to the dense union.  Contractivity of
conditional expectations extends the equality and convergence by
continuity.
\end{proof}

\begin{corollary}[$\Z_4$ fusion and expectation]\label{cor:Z4}
For $G=\Z_4$, the limiting homogeneous subspaces
$\B_k=\overline{Au^k}$ satisfy
\[
 \B_k\B_\ell\subset\B_{k+\ell},\qquad
 \B_k^*=\B_{-k},\qquad u^4=\id.
\tag{6.7}\label{eq:grades}
\]
Thus the grade-one field obeys
$\Psi_\lambda^4\in\A$ and
$\Psi_\lambda^2\in\B_2$, identified by
\eqref{eq:lambda-order} with the vector/parity class.  The limit
expectation is
\[
 E(x)=\frac14\sum_{r=0}^3\alpha_r(x),\qquad E(x)\ge\frac14x
\quad(x\ge0).
\tag{6.8}\label{eq:quarterE}
\]
\end{corollary}

\begin{proof}
For $a,b\in A$,
$(au^k)(bu^\ell)=a\alpha_k(b)u^{k+\ell}$ and
$(au^k)^*=\alpha_{-k}(a^*)u^{-k}$.  The remaining statements follow from
Theorem~\ref{thm:crossedlimit}.
\end{proof}

\begin{remark}[Inner finite shadows versus outer type-III actions]
At a finite matrix stage the clock action can be inner, in which case
$A_N\rtimes\Z_4\cong A_N\otimes C^*(\Z_4)$ and the inclusion need not
model an irreducible outer fixed-point subfactor.  The limit theorem
above remains valid, but index \emph{equality} and Q-system
irreducibility require outerness/saturation of the limiting local action.
Those properties belong to L2/\textup{(ALG-EXH)}, not to the formal
crossed-product exchange.
\end{remark}

\section{L4: holomorphy and identification with \texorpdfstring{$(E_8)_1$}{E8 level 1}}
\label{sec:L4}

